% Generated draft: requires mathematical and source review.
\documentclass[11pt]{article}
\usepackage[margin=25mm]{geometry}
\usepackage{fontspec}
\setmainfont{Latin Modern Roman}
\usepackage{amsmath,amssymb,mathtools,unicode-math}
\setmathfont{Latin Modern Math}
\usepackage{enumitem,needspace,xurl,hyperref}
\setlength{\parindent}{0pt}
\setlength{\parskip}{5pt}
\setlength{\emergencystretch}{3em}
\newcommand{\N}{\mathbb{N}}
\newcommand{\Z}{\mathbb{Z}}
\newcommand{\Q}{\mathbb{Q}}
\newcommand{\R}{\mathbb{R}}
\newcommand{\C}{\mathbb{C}}
\newcommand{\F}{\mathbb{F}}
\newcommand{\A}{\mathbb{A}}
\newcommand{\PP}{\mathbb{P}}
\newcommand{\E}{\mathbb{E}}
\newcommand{\eps}{\varepsilon}
\newcommand{\dd}{\,\mathrm{d}}
\begin{document}
\section*{UnsolvedMath ID: 30002867}
\textbf{Source problem code:} catalogue entry
\subsection{Definitions and mathematical statement}
Let \(d\geq 1\) and \(R=\C[x_1,\ldots,x_d]\). An ideal projector is a \(\C\)-linear map \(P:R\to R\) such that \(P^2=P\) and \(I:=\ker P\) is an ideal. Put \(A=R/I\). For each \(j\in\{1,\ldots,d\}\), define the multiplication operator
\[
\mu_j:A\longrightarrow A,\qquad \mu_j(f+I)=x_jf+I.
\]
These operators commute. In the finite-dimensional, nonzero setting \(0<N:=\dim_{\C}A<\infty\), choose a common basis of \(A\) and let \(M_j\in\operatorname{Mat}_N(\C)\) represent \(\mu_j\). A change of basis simultaneously conjugates all the matrices.

For a precise global interpretation, call a proper ideal \(I\subset R\) a complete intersection if \(I=(f_1,\ldots,f_r)\) for an \(R\)-regular sequence: for each \(k\), multiplication by the class of \(f_k\) is injective on \(R/(f_1,\ldots,f_{k-1})\). In the finite-dimensional, nonzero setting above, such a sequence has length \(d\).

The requested characterization is a necessary and sufficient condition, expressed in terms of the associated commuting matrices \((M_1,\ldots,M_d)\) and invariant under simultaneous change of basis, for \(\ker P\) to be a complete intersection.

The supplied statement does not explicitly impose finite rank or exclude \(A=0\); the finite-dimensional, nonzero setting above makes the matrix formulation precise but requires confirmation against the source. The global regular-sequence meaning of complete intersection is also an explicit interpretation requiring confirmation; a condition imposed only after localization would be a different target. The original question refers to matrices defined by equation \((1)\), whose text is absent from the supplied record. The multiplication-operator formulation here follows the supplied restored statement.
\subsection{Short English statement}
Characterize when the kernel of an ideal projector on \(\C[x_1,\ldots,x_d]\) is a complete intersection using the commuting matrices for multiplication by the coordinate variables on the quotient. The finite-rank setting, the zero-quotient boundary case, and the intended global or local meaning of complete intersection require source confirmation.
\subsection{Sources}
\begin{itemize}
\item Local UnsolvedMath registry, permanent ID \(30002867\), \emph{Matrix Characterization of Complete Intersections}; exact source problem code \texttt{catalogue entry}, in the Oberwolfach Reports Open Problems collection. No individual proposer is supplied; the proposed year is recorded as 2015.
\item Supplied report citation: \emph{Multivariate Splines and Algebraic Geometry (2016)}, Oberwolfach Reports, pp.\ 1139--1200, DOI \texttt{10.4171/owr/2015/21}. Supplied literature-source label: Oberwolfach Report 21. \url{https://doi.org/10.4171/owr/2015/21}. The citation's parenthetical year 2016 differs from the proposed year 2015 and the year appearing in the DOI; this discrepancy is retained without resolution.
\item Shekhtman, \emph{Some Tidbits on Ideal Projectors}. \url{https://eudml.org/doc/229343}. No further bibliographic details are supplied.
\item Literature assessment checked at \texttt{2026-08-21}: the supplied curation labels the problem open and states that the cited projector theory provides criteria and examples but no general characterization. This is supplied dated curation, not an independently verified literature assessment.
\item The registry labels the statement \texttt{corrected\_verified} and reports that the official Oberwolfach passage was checked directly and its definitions and notation restored without strengthening the question. This verification claim is supplied provenance and has not been independently checked in preparing this draft.
\end{itemize}
\end{document}
